% Generated draft: requires mathematical and source review.
\documentclass[11pt]{article}
\usepackage[margin=25mm]{geometry}
\usepackage{fontspec}
\setmainfont{Latin Modern Roman}
\usepackage{amsmath,amssymb,mathtools,unicode-math}
\setmathfont{Latin Modern Math}
\usepackage{enumitem,needspace,xurl,hyperref}
\setlength{\parindent}{0pt}
\setlength{\parskip}{5pt}
\setlength{\emergencystretch}{3em}
\newcommand{\N}{\mathbb{N}}
\newcommand{\Z}{\mathbb{Z}}
\newcommand{\Q}{\mathbb{Q}}
\newcommand{\R}{\mathbb{R}}
\newcommand{\C}{\mathbb{C}}
\newcommand{\F}{\mathbb{F}}
\newcommand{\A}{\mathbb{A}}
\newcommand{\PP}{\mathbb{P}}
\newcommand{\E}{\mathbb{E}}
\newcommand{\eps}{\varepsilon}
\newcommand{\dd}{\,\mathrm{d}}
\begin{document}
\section*{UnsolvedMath ID: 800124}
\textbf{Source problem code:} catalogue entry
\subsection{Definitions and mathematical statement}
Fix a dimension \(q\); the supplied source does not specify its intended range. Let \(\Gamma_q\) be the groupoid whose objects are points of \(\R^q\) and whose arrows are germs \([f]_x\) of local diffeomorphisms, with source \(x\), target \(f(x)\), and composition induced by composition of representatives. Let
\[
G_q=\{[f]_0\in\Gamma_q:f(0)=0\}
\]
be its isotropy group at the origin. Write \(B\) for classifying space. The normal bundle of the universal \(G_q\)-structure is classified by a map
\[
\nu_G:BG_q\longrightarrow BGL_q(\R),
\]
induced by \([f]_0\mapsto D_0f\). Define
\[
B\bar G_q:=\operatorname{hofib}_{*}(\nu_G),
\]
where \(*\) is the basepoint corresponding to the trivial rank-\(q\) vector bundle.

The supplied statement uses \(B\bar\Gamma_q\) without defining it. Under the conventional interpretation, let \(\nu_\Gamma:B\Gamma_q\to BGL_q(\R)\) classify the normal bundle of the universal \(\Gamma_q\)-structure and put
\[
B\bar\Gamma_q:=\operatorname{hofib}_{*}(\nu_\Gamma).
\]
The inclusion \(G_q\hookrightarrow\Gamma_q\), together with its compatibility with the normal-bundle maps, induces a map
\[
\iota:B\bar G_q\longrightarrow B\bar\Gamma_q.
\]
Interpreting the source phrase ``homotopically trivial in'' as saying that this map is null-homotopic, the question is whether there exist a point \(c\in B\bar\Gamma_q\) and a continuous map
\[
H:B\bar G_q\times[0,1]\longrightarrow B\bar\Gamma_q
\]
such that, for every \(x\in B\bar G_q\),
\[
H(x,0)=\iota(x),\qquad H(x,1)=c.
\]
This precise formulation is conditional on the stated interpretation of the target and of ``homotopically trivial.'' The source does not explicitly specify the differentiability class, the topologies used for the germ groupoid and its isotropy group, the range of \(q\), or whether a based null-homotopy is required. These conventions require confirmation.
\subsection{Short English statement}
Question 65 asks whether the natural map from the normal-bundle homotopy fibre for germs fixing the origin to the corresponding fibre for all local diffeomorphism germs is null-homotopic. The supplied wording leaves some conventions and the exact meaning of ``homotopically trivial'' unresolved.
\subsection{Sources}
\begin{itemize}
\item Local UnsolvedMath registry, permanent ID \(800124\), exact source problem code \texttt{catalogue entry}. The supplied record attributes the question to Nariman and Yazdi (2025) and labels it open; that label is supplied source data, not an independently verified status assessment.
\item Nariman and Yazdi, \emph{Problem List on Foliations and Diffeomorphism Groups} (2025), Question 65. Supplied source URL: \url{https://arxiv.org/abs/2503.18274}. The record reports access on 2026-10-16 and extraction by \texttt{arxiv\_latex}; no independent source visit or verification is asserted here.
\item Supplied attribution/context: ``Foliations with only one leaf'': see Tsuboi, reference label \texttt{tsuboi1984gamma1}. The supplied record gives no full bibliographic entry or URL for this reference.
\end{itemize}
\end{document}
